\SourceLeaf{132}
Quia
\[
(\alpha\alpha'+\beta\gamma')(\gamma\beta'+\delta\delta')
 -(\alpha\beta'+\beta\delta')(\gamma\alpha'+\delta\gamma')
 =(\alpha\delta-\beta\gamma)(\alpha'\delta'-\beta'\gamma')
\]
adeoque positivus, si tum \(\alpha\delta-\beta\gamma\) tum \(\alpha'\delta'-\beta'\gamma'\) positivus aut uterque negativus, negativus vero si alter horum numerorum positivus alter negativus: forma \(F\) formam \(F''\) \textit{proprie} implicabit, si \(F\) ipsam \(F'\) et \(F'\) ipsam \(F''\) eodem modo implicant, \textit{improprie} si diverso.

Hinc sequitur, si quotcunque formae habeantur \(F,F',F'',F'''\) etc., quarum quaevis sequentem implicet, primam implicaturam esse ultimam, et quidem \textit{proprie}, si multitudo formarum, quae sequentem suam improprie implicant, fuerit par, \textit{improprie} si multitudo haec impar.

\textit{Si forma \(F\) formae \(F'\) est aequivalens, formaque \(F'\) formae \(F''\): forma \(F\) formae \(F''\) aequivalens erit, et quidem proprie, si forma \(F\) formae \(F'\) eodem modo aequivalet ut forma \(F'\) formae \(F''\), improprie, si diverso.}

Quia enim formae \(F,F'\), his \(F',F''\), respective sunt aequivalentes, tum illae has resp. implicabunt, adeoque \(F\) ipsam \(F''\), tum hae illas. Quare \(F,F''\) aequivalentes erunt. Ex praec. vero sequitur, \(F\) ipsam \(F''\) proprie vel improprie implicare, prout \(F\) ipsi \(F'\) et \(F'\) ipsi \(F''\) eodem modo vel diverso sint aequivalentes, ut et \(F''\) ipsam \(F\): quare in priori casu \(F,F''\) proprie, in posteriori improprie aequivalentes erunt.

\textit{Formae \((a,-b,c)\), \((c,b,a)\), \((c,-b,a)\) formae \((a,b,c)\) aequivalent, et quidem duae priores improprie, ultima proprie.}

Nam \(axx+2bxy+cyy\) transit in \(ax'x'-2bx'y'+cy'y'\), ponendo \(x=x'+0.y'\), \(y=0.x'-y'\), quae transformatio est impropria propter \(1\times-1-0.0=-1\); in formam \(cx'x'+2bx'y'+ay'y'\) vero per transformationem impropriam \(x=0.x'+y'\), \(y=x'+0.y'\); et in formam \(cx'x'-2bx'y'+ay'y'\) per propriam \(x=0.x'-y'\), \(y=x'+0.y'\).

Hinc manifestum est, quamvis formam, formae \((a,b,c)\) aequivalentem, vel ipsi, vel formae \((a,-b,c)\) \textit{proprie} aequivalere; similiterque, si quae forma formam \((a,b,c)\) implicet aut sub ipsa contineatur, eam vel formam \((a,b,c)\) vel formam \((a,-b,c)\) \textit{proprie} implicare, aut sub alterutra \textit{proprie} contineri. Formas \((a,b,c)\), \((a,-b,c)\) \textit{oppositas} vocabimus.

\SourceLeaf{133}
\begin{center}
\textit{Formae contiguae.}
\end{center}
\Article{160}
Si formae \((a,b,c)\), \((a',b',c')\) eundem determinantem habent, insuperque est \(c=a'\) et \(b\equiv-b'\; (\mathrm{mod.}\,c)\), sive \(b+b'\equiv0\; (\mathrm{mod.}\,c)\), formas has \textit{contiguas} dicemus, et quidem, quando determinatione accuratiori opus est, priorem posteriori \textit{a parte prima}, posteriorem priori \textit{a parte ultima} contiguam dicemus.

Ita \textit{ex. gr.} forma \((7,3,2)\) formae \((3,4,7)\) a parte ultima contigua, forma \((3.1.3)\) oppositae suae \((3,-1,3)\) ab utraque parte.

\textit{Formae contiguae semper sunt proprie aequivalentes.} Nam forma \(axx+2bxy+cyy\) transit in formam contiguam \(cx'x'+2b'x'y'+c'y'y'\) per substitutionem \(x=-y'\), \(y=x'+\frac{b+b'}{c}y'\) (quae est propria ob \(0\times\frac{b+b'}{c}-1\times-1=1\)), uti per evolutionem adiumento aequationis \(bb-ac=b'b'-cc'\) facile probatur; \(\frac{b+b'}{c}\) vero per hyp. est integer. --- Ceterum hae definitiones et conclusiones locum non habent. si \(c=a'=0\). Hic vero casus occurrere nequit, nisi in formis quarum determinans est numerus quadratus.

Formae \((a,b,c)\), \((a',b',c')\) proprie aequivalentes sunt, si \(a=a'.\ b\equiv b'\; (\mathrm{mod.}\,a)\). Forma enim \((a,b,c)\) formae \((c,-b,a)\) proprie aequivalet (art. praec.), haec vero formae \((a',b',c')\) a parte prima contigua erit.

\begin{center}
\textit{Divisores communes co\"efficientium formarum.}
\end{center}
\Article{161}
\textit{Si forma \((a,b,c)\) formam \((a',b',c')\) implicat, quivis divisor communis numerorum \(a,b,c\) etiam numeros \(a',b',c'\) metietur, et quivis divisor communis numerorum \(a,2b,c\) ipsos \(a',2b',c'\).}

Si enim forma \(axx+2bxy+cyy\) per substitutiones \(x=\alpha x'+\beta y'\), \(y=\gamma x'+\delta y'\) in formam \(a'x'x'+2b'x'y'+c'y'y'\) transit: habebuntur hae aequationes
\[
\begin{aligned}
a\alpha\alpha+2b\alpha\gamma+c\gamma\gamma&=a'\\
a\alpha\beta+b(\alpha\delta+\beta\gamma)+c\gamma\delta&=b'\\
a\beta\beta+2b\beta\delta+c\delta\delta&=c'
\end{aligned}
\]
unde propositio statim sequitur (pro parte secunda propos. loco aequationis secundae hanc adhibendo \(2a\alpha\beta+2b(\alpha\delta+\beta\gamma)+2c\gamma\delta=2b'\)).

\SourceLeaf{134}
Hinc sequitur maximum divisorem communem numerorum \(a,b(2b),c\) simul metiri divisorem communem maximum numerorum \(a',b'(2b'),c'\). Quodsi igitur insuper forma \((a',b',c')\) formam \((a,b,c)\) implicat, \textit{i. e.} formae sunt aequivalentes, divisor communis maximus numerorum \(a,b(2b),c\), divisori communi maximo numerorum \(a',b'(2b'),c'\) aequalis erit, quoniam tum ille hunc metiri debet, tum hic illum. Si itaque in hoc casu \(a,b(2b),c\) divisorem communem non habent, \textit{i. e.} si maximus \(=1\), etiam \(a',b'(2b'),c'\) divisorem communem non habebunt.

\Needspace{18\baselineskip}
\begin{center}
\textit{Nexus omnium transformationum similium formae datae in formam datam.}
\end{center}
\Article{162}
\textsc{Problema.} \textit{Si forma}
\[
AXX+2BXY+CYY\quad\mathinner{.\,.\,.}\quad F
\]
\textit{formam}
\[
axx+2bxy+cyy\quad\mathinner{.\,.\,.\,.\,.}\quad f
\]
\textit{implicat, atque transformatio aliqua illius in hanc est data: ex hac omnes reliquas transformationes ipsi similes deducere.}

\textit{Solutio.} Sit transformatio data haec \(X=\alpha x+\beta y\), \(Y=\gamma x+\delta y\), ponamusque primo aliam huic similem datam esse \(X=\alpha'x+\beta'y\), \(Y=\gamma'x+\delta'y\), ut quid inde sequatur investigemus. Tum positis determinantibus formarum \(F,f,\ =D,d\), atque \(\alpha\delta-\beta\gamma=e\), \(\alpha'\delta'-\beta'\gamma'=e'\), erit (art. 157), \(d=Dee=De'e'\), et quum ex hyp. \(e,e'\) eadem signa habeant, \(e=e'\). Habebuntur autem sequentes sex aequationes:
\[
\begin{aligned}
A\alpha\alpha+2B\alpha\gamma+C\gamma\gamma&=a&\quad\mathinner{.\,.\,.\,.\,.\,.\,.}\quad&[1]\\
A\alpha'\alpha'+2B\alpha'\gamma'+C\gamma'\gamma'&=a&\quad\mathinner{.\,.\,.\,.\,.\,.\,.}\quad&[2]\\
A\alpha\beta+B(\alpha\delta+\beta\gamma)+C\gamma\delta&=b&\quad\mathinner{.\,.\,.\,.\,.\,.}\quad&[3]\\
A\alpha'\beta'+B(\alpha'\delta'+\beta'\gamma')+C\gamma'\delta'&=b&\quad\mathinner{.\,.\,.\,.\,.\,.}\quad&[4]\\
A\beta\beta+2B\beta\delta+C\delta\delta&=c&\quad\mathinner{.\,.\,.\,.\,.\,.\,.}\quad&[5]\\
A\beta'\beta'+2B\beta'\delta'+C\delta'\delta'&=c&\quad\mathinner{.\,.\,.\,.\,.\,.\,.}\quad&[6]
\end{aligned}
\]
Si brevitatis gratia numeros
\[
\begin{gathered}
A\alpha\alpha'+B(\alpha\gamma'+\gamma\alpha')+C\gamma\gamma'\\
A(\alpha\beta'+\beta\alpha')+B(\alpha\delta'+\beta\gamma'+\gamma\beta'+\delta\alpha')+C(\gamma\delta'+\delta\gamma')\\
A\beta\beta'+B(\beta\delta'+\delta\beta')+C\delta\delta'
\end{gathered}
\]
\SourceLeaf{135}
per \(a',2b',c'\) designamus, ex aequ. praec. sequentes novas deducemus\HistoricalNote{*)}{Origo harum aequationum haec est: \(7\) fit ex \(1.2\) (\textit{i. e.} si aequatio (1) in aequationem (2) multiplicatur, sive potius, si illius pars prior in partem priorem huius multiplicatur, illiusque pars posterior in posteriorem huius, productaque aequalia ponuntur); \(8\) ex \(1.4+2.3\); sequens quae non est numerata ex \(1.6+2.5+3.4+3.4\); sequens non numerata ex \(3.4\); \(11\) ex \(3.6+4.5\); \(12\) ex \(5.6\). Simili designatione etiam in sequentibus semper utemur. Evolutionem vero lectoribus relinquere debemus.}:
\[
\begin{aligned}
a'a'-D(\alpha\gamma'-\gamma\alpha')^2&=aa&\quad\mathinner{.\,.\,.\,.\,.\,.\,.\,.}\quad&[7]\\
2a'b'-D(\alpha\gamma'-\gamma\alpha')(\alpha\delta'+\beta\gamma'-\gamma\beta'-\delta\alpha')&=2ab&\quad\mathinner{.\,.}\quad&[8]\\
4b'b'-D\bigl((\alpha\delta'+\beta\gamma'-\gamma\beta'-\delta\alpha')^2+2ee'\bigr)&=2bb+2ac
\end{aligned}
\]
unde fit, addendo \(2Dee'=2d=2bb-2ac\),
\[
\begin{aligned}
4b'b'-D(\alpha\delta'+\beta\gamma'-\gamma\beta'-\delta\alpha')^2&=4bb&\quad\mathinner{.\,.\,.\,.\,.}\quad&[9]\\
a'c'-D(\alpha\delta'-\gamma\beta')(\beta\gamma'-\delta\alpha')&=bb
\end{aligned}
\]
unde subtrahendo \(D(\alpha\delta-\beta\gamma)(\alpha'\delta'-\beta'\gamma')=bb-ac\) fit
\[
\begin{aligned}
a'c'-D(\alpha\gamma'-\gamma\alpha')(\beta\delta'-\delta\beta')&=ac&\quad\mathinner{.\,.\,.\,.\,.\,.}\quad&[10]\\
2b'c'-D(\alpha\delta'+\beta\gamma'-\gamma\beta'-\delta\alpha')(\beta\delta'-\delta\beta')&=2bc&\quad\mathinner{.\,.}\quad&[11]\\
c'c'-D(\beta\delta'-\delta\beta')^2&=cc&\quad\mathinner{.\,.\,.\,.\,.\,.\,.\,.}\quad&[12]
\end{aligned}
\]
Ponamus iam, divisorem communem maximum numerorum \(a,2b,c\) esse \(m\) numerosque \(\mathfrak{A},\mathfrak{B},\mathfrak{C}\) ita determinatos, ut fiat
\[
\mathfrak{A}a+2\mathfrak{B}b+\mathfrak{C}c=m
\]
(art. 40); multiplicentur aequationes \(7,8,9,10,11,12\) resp. per \(\mathfrak{A}\mathfrak{A},\allowbreak 2\mathfrak{A}\mathfrak{B},\allowbreak \mathfrak{B}\mathfrak{B},\allowbreak 2\mathfrak{A}\mathfrak{C},\allowbreak 2\mathfrak{B}\mathfrak{C},\allowbreak \mathfrak{C}\mathfrak{C}\) summenturque producta. Quodsi iam brevitatis caussa ponimus
\[
\begin{aligned}
\mathfrak{A}a'+2\mathfrak{B}b'+\mathfrak{C}c'&=T&\quad\mathinner{.\,.\,.\,.\,.\,.\,.\,.}\quad&[13]\\
\mathfrak{A}(\alpha\gamma'-\gamma\alpha')&+\mathfrak{B}(\alpha\delta'+\beta\gamma'-\gamma\beta'-\delta\alpha')&&\\
{}+\mathfrak{C}(\beta\delta'-\delta\beta')&=U&\quad\mathinner{.\,.}\quad&[14]
\end{aligned}
\]
ubi \(T,U\) manifesto erunt integri, prodibit
\[
TT-DUU=mm
\]
Deducti itaque sumus ad hanc conclusionem elegantem, \textit{ex binis quibuscunque transformationibus similibus formae \(F\) in \(f\) sequi solutionem aequationis indeterminatae \(tt-Duu=mm\), in integris}, scilicet \(t=T,\ u=U\). Ceterum quum in
\SourceLeaf{136}
ratiociniis nostris non supposuerimus, transformationes esse \textit{diversas}: \textit{una} adeo transformatio bis considerata solutionem praebere debet. Tum vero fit propter \(\alpha'=\alpha,\ \beta'=\beta\) etc. \(a'=a,\ b'=b,\ c'=c\), adeoque \(T=m,\ U=0\), quae solutio per se est obvia.

Iam primam transformationem solutionemque aequationis indeterminatae tamquam cognitas consideremus, et quomodo hinc altera transformatio deduci possit, sive quomodo \(\alpha',\beta',\gamma',\delta'\), ab his \(\alpha,\beta,\gamma,\delta,\ T,U\) pendeant, investigemus. Ad hunc finem multiplicamus primo aequationem \([1]\) per \(\delta\alpha'-\beta\gamma'\), \([2]\) per \(\alpha\delta'-\gamma\beta'\), \([3]\) per \(\alpha\gamma'-\gamma\alpha'\), \([4]\) per \(\gamma\alpha'-\alpha\gamma'\), addimusque producta, unde prodibit
\[
(e+e')a'=(\alpha\delta'-\beta\gamma'-\gamma\beta'+\delta\alpha')a\quad\mathinner{.\,.\,.\,.\,.\,.}\quad[15]
\]
Simili modo fit ex
\[
(\delta\beta'-\beta\delta')([1]-[2])+(\alpha\delta'-\beta\gamma'-\gamma\beta'+\delta\alpha')([3]+[4])+(\alpha\gamma'-\gamma\alpha')([5]-[6])
\]
\[
2(e+e')b'=2(\alpha\delta'-\beta\gamma'-\gamma\beta'+\delta\alpha')b\quad\mathinner{.\,.\,.\,.\,.}\quad[16]
\]
Denique ex \((\delta\beta'-\beta\delta')([3]-[4])+(\alpha\delta'-\gamma\beta')[5]+(\delta\alpha'-\beta\gamma')[6]\) prodit:
\[
(e+e')c'=(\alpha\delta'-\beta\gamma'-\gamma\beta'+\delta\alpha')c\quad\mathinner{.\,.\,.\,.\,.\,.}\quad[17]
\]
Substituendo hos valores \((15,16,17)\) in \(13\) fit
\[
(e+e')T=(\alpha\delta'-\beta\gamma'-\gamma\beta'+\delta\alpha')(\mathfrak{A}a+2\mathfrak{B}b+\mathfrak{C}c)
\]
sive
\[
2eT=(\alpha\delta'-\beta\gamma'-\gamma\beta'+\delta\alpha')m\quad\mathinner{.\,.\,.\,.\,.\,.\,.}\quad[18]
\]
unde \(T\) multo facilius deduci potest, quam ex \([13]\). --- Combinando hanc aequationem cum \(15,16,17\) obtinetur \(ma'=Ta,\ 2mb'=2Tb,\ mc'=Tc\). Quos valores ipsorum \(a',2b',c'\) in aequ. \(7\text{---}12\) substituendo et loco ipsius \(TT\) scribendo \(mm+DUU\), transeunt illae post mutationes debitas in has
\[
\begin{aligned}
(\alpha\gamma'-\gamma\alpha')^2mm&=aaUU\\
(\alpha\gamma'-\gamma\alpha')(\alpha\delta'+\beta\gamma'-\gamma\beta'-\delta\alpha')mm&=2abUU
\end{aligned}
\]
\SourceLeaf{137}
\[
\begin{aligned}
(\alpha\delta'+\beta\gamma'-\gamma\beta'-\delta\alpha')^2mm&=4bbUU\\
(\alpha\gamma'-\gamma\alpha')(\beta\delta'-\delta\beta')mm&=acUU\\
(\alpha\delta'+\beta\gamma'-\gamma\beta'-\delta\alpha')(\beta\delta'-\delta\beta')mm&=2bcUU\\
(\beta\delta'-\delta\beta')^2mm&=ccUU
\end{aligned}
\]
Hinc adiumento aequationis \([14]\) et huius \(\mathfrak{A}a+2\mathfrak{B}b+\mathfrak{C}c=m\), facile deducitur (multiplicando primam, secundam, quartam; secundam, tertiam, quintam; quartam, quintam, sextam, resp. per \(\mathfrak{A},\mathfrak{B},\mathfrak{C}\) addendoque producta):
\[
\begin{aligned}
(\alpha\gamma'-\gamma\alpha')Umm&=maUU\\
(\alpha\delta'+\beta\gamma'-\gamma\beta'-\delta\alpha')Umm&=2mbUU\\
(\beta\delta'-\delta\beta')Umm&=mcUU
\end{aligned}
\]
atque hinc, dividendo per \(mU\)\HistoricalNote{*)}{Hoc non liceret, si esset \(U=0\): tunc vero aequationum \(19,20,21\) veritas statim ex prima, tertia et sexta praecedentium sequeretur.}
\[
\begin{aligned}
aU&=(\alpha\gamma'-\gamma\alpha')m&\quad\mathinner{.\,.\,.\,.\,.\,.\,.\,.\,.}\quad&[19]\\
2bU&=(\alpha\delta'+\beta\gamma'-\gamma\beta'-\delta\alpha')m&\quad\mathinner{.\,.\,.\,.\,.}\quad&[20]\\
cU&=(\beta\delta'-\delta\beta')m&\quad\mathinner{.\,.\,.\,.\,.\,.\,.\,.\,.}\quad&[21]
\end{aligned}
\]
ex quarum aequationum aliqua \(U\) multo facilius quam ex \([14]\) deduci potest. --- Simul hinc colligitur, quomodocunque \(\mathfrak{A},\mathfrak{B},\mathfrak{C}\) determinentur (quod infinitis modis diversis fieri potest), tum \(T\) tum \(U\) eundem valorem adipisci.

Iam si aequatio \(18\) multiplicatur per \(\alpha\), \(19\) per \(2\beta\), \(20\) per \(-\alpha\), fit per additionem
\[
2\alpha eT+2(\beta a-\alpha b)U=2(\alpha\delta-\beta\gamma)\alpha'm=2e\alpha'm
\]
Simili modo fit ex \(\beta[18]+\beta[20]-2\alpha[21]\)
\[
2\beta eT+2(\beta b-\alpha c)U=2(\alpha\delta-\beta\gamma)\beta'm=2e\beta'm
\]
Porro ex \(\gamma[18]+2\delta[19]-\gamma[20]\) fit
\[
2\gamma eT+2(\delta a-\gamma b)U=2(\alpha\delta-\beta\gamma)\gamma'm=2e\gamma'm
\]
Tandem ex \(\delta[18]+\delta[20]-2\gamma[21]\) prodit
\[
2\delta eT+2(\delta b-\gamma c)U=2(\alpha\delta-\beta\gamma)\delta'm=2e\delta'm
\]

\SourceLeaf{138}
In quibus formulis, si pro \(a,b,c\) valores ex \(1,3,5\) substituantur, fit
\[
\begin{aligned}
\alpha'm&=\alpha T-(\alpha B+\gamma C)U\\
\beta'm&=\beta T-(\beta B+\delta C)U\\
\gamma'm&=\gamma T+(\alpha A+\gamma B)U\\
\delta'm&=\delta T+(\beta A+\delta B)U\text{\textsuperscript{*)}}
\end{aligned}
\]
\HistoricalDisplayNote{*)}{Hinc facile deducitur
\[
\begin{aligned}
AeU&=(\delta\gamma'-\gamma\delta')m\\
2BeU&=(\alpha\delta'-\delta\alpha'+\gamma\beta'-\beta\gamma')m\\
CeU&=(\beta\alpha'-\alpha\beta')m
\end{aligned}
\]}

Ex analysi praec. sequitur, nullam transformationem formae \(F\) in \(f\) propositae similem dari, quae non sit contenta sub formula
\[
\begin{aligned}
X&=\tfrac{1}{m}(\alpha t-(\alpha B+\gamma C)u)x+\tfrac{1}{m}(\beta t-(\beta B+\delta C)u)y\\
Y&=\tfrac{1}{m}(\gamma t+(\alpha A+\gamma B)u)x+\tfrac{1}{m}(\delta t+(\beta A+\delta B)u)y\quad\mathinner{.\,.\,.}\quad(\mathrm{I})
\end{aligned}
\]
designantibus \(t,u\) indefinite omnes numeros integros aequationi \(tt-Duu=mm\) satisfacientes. Hinc vero concludere nondum possumus, omnes valores ipsorum \(t,u\), aequationi illi satisfacientes, in formula \((\mathrm{I})\) substitutos, transformationes idoneas praebere. At

1. Formam \(F\) per substitutionem, e quibusvis ipsorum \(t,u\) valoribus ortam, semper in formam \(f\) transmutari, per evolutionem confirmari facile potest adiumento aequationum \(1,3,5\) et huius \(tt-Duu=mm\). Calculum prolixiorem quam difficiliorem brevitatis gratia supprimimus.

2. Quaevis transformatio ex formula deducta propositae erit similis. Namque
\[
\begin{aligned}
&\tfrac{1}{m}(\alpha t-(\alpha B+\gamma C)u)\times\tfrac{1}{m}(\delta t+(\beta A+\delta B)u)\\
&\quad-\tfrac{1}{m}(\beta t-(\beta B+\delta C)u)\times\tfrac{1}{m}(\gamma t+(\alpha A+\gamma B)u)\\
&\qquad=\tfrac{1}{mm}(\alpha\delta-\beta\gamma)(tt-Duu)=\alpha\delta-\beta\gamma
\end{aligned}
\]
3. Si formae \(F,f\) determinantes inaequales habent, fieri potest, ut formula \((\mathrm{I})\) pro quibusdam valoribus ipsorum \(t,u\) praebeat substitutiones, quae \textit{fractiones} implicent. adeoque reiici debeant. Omnes vero reliquae erunt transformationes idoneae, aliaeque praeter ipsas non dabuntur.

4. Si vero formae \(F,f\) eundem determinantem habent adeoque sunt \textit{aequivalentes}, formula \((\mathrm{I})\) nullas transformationes quae fractiones implicent praebe\SourceLeaf{139}bit, adeoque in hoc casu solutionem completam problematis exhibebit. Illud vero ita demonstramus.

Ex theoremate art. praec. sequitur in hocce casu, \(m\) simul fore divisorem communem numerorum \(A,2B,C\). Quoniam \(tt-Duu=mm\), fit \(tt-BBuu=mm-ACuu\), quare \(tt-BBuu\) per \(mm\) divisibilis erit: hinc etiam a potiori \(4tt-4BBuu\) adeoque (quia \(2B\) per \(m\) divisibilis) etiam \(4tt\) per \(mm\) et proin \(2t\) per \(m\). Hinc \(\frac{2}{m}(t+Bu),\ \frac{2}{m}(t-Bu)\) erunt integri, et quidem (quoniam differentia inter ipsos \(\frac{4}{m}Bu\) est par) aut uterque par aut uterque impar. Si uterque impar esset. etiam productum impar foret, quod tamquam quadruplum numeri \(\frac{1}{mm}(tt-BBuu)\), quem integrum esse modo ostendimus, necessario par: quare hic casus est impossibilis, adeoque \(\frac{2}{m}(t+Bu),\ \frac{2}{m}(t-Bu)\) semper pares, unde \(\frac{1}{m}(t+Bu),\ \frac{1}{m}(t-Bu)\) erunt integri. Hinc vero nullo negotio deducitur, omnes quatuor coefficientes in \((\mathrm{I})\) semper esse integros. \mbox{\textit{Q. E. D.}}

Ex praecedentibus colligitur, si omnes solutiones aequationis \(tt-Duu=mm\) habeantur, omnes transformationes formae \((A,B,C)\) in \((a,b,c)\) transf. datae similes inde derivari. Illas vero in sequentibus invenire docebimus. Hic tantummodo observamus multitudinem solutionum semper esse finitam, quando \(D\) sit negativus. aut positivus simulque quadratus: quando vero \(D\) positivus non quadratus, infinitam. Quando hic casus locum habet, simulque \(D\) non \(=d\) (supra \(3^0\)), disquiri insuper deberet, quomodo ii valores ipsorum \(t,u\), qui substitutiones a fractionibus liberas, ab iis, qui fractas producunt, a priori dignosci possint. Sed pro hocce casu infra aliam methodum ab hoc incommodo liberam exponemus (art. 214).

\textit{Exempl.} Forma \(xx+2yy\) per substitutionem propriam \(x=2x'+7y'\), \(y=x'+5y'\) transit in formam \((6,24,99)\): desiderantur \textit{omnes} transformationes propriae formae illius in hanc. Hic \(D=-2,\ m=3\), adeoque aequatio solvenda haec: \(tt+2uu=9\). Huic sex modis diversis satisfit ponendo scilicet \(t=3,-3,1,-1,1,-1;\ u=0,0,2,2,-2,-2\) resp. Solutio tertia et sexta dant substitutiones in fractis, adeoque sunt reiiciendae; ex reliquis sequuntur quatuor substitutiones:
\[
x=\left\{\begin{aligned}
&2x'+7y'\\
&-2x'-7y'\\
&-2x'-9y'\\
&2x'+9y'
\end{aligned}\right.
\qquad
y=\left\{\begin{aligned}
&x'+5y'\\
&-x'-5y'\\
&x'+3y'\\
&-x'-3y'
\end{aligned}\right.
\]
(quarum prima est proposita).

\SourceLeaf{140}
\Needspace{8\baselineskip}
\begin{center}
\textit{Formae ancipites.}
\end{center}
\Article{163}
Iam supra obiter diximus fieri posse ut forma aliqua, \(F\), aliam, \(F'\), tam proprie quam improprie implicet. Perspicuum est hoc evenire, si inter formas \(F,F'\) alia \(G\) interponi possit, ita ut \(F\) ipsam \(G\), \(G\) ipsam \(F'\) implicet, formaque \(G\) ita sit comparata, ut sibi ipsa sit improprie aequivalens. Si enim \(F\) ipsam \(G\) proprie vel improprie implicare supponitur: quum \(G\) ipsam \(G\) improprie implicet, \(F\) ipsam \(G\) improprie vel proprie (resp.) implicabit, adeoque, in utroque casu, tam proprie quam improprie (art. 159). Eodem modo hinc deducitur, quomodocunque \(G\) ipsam \(F'\) implicare supponatur, \(F\) semper ipsam \(F'\) tum proprie tum improprie implicare debere. --- Tales vero formas dari, quae sibi ipsae sint improprie aequivalentes, videtur in casu maxime obvio, ubi formae terminus medius \(=0\). Talis enim forma sibi ipsa erit opposita (art. 159) adeoque improprie aequivalens. Generalius quaevis forma \((a,b,c)\) hac proprietate est praedita, in qua \(2b\) per \(a\) est divisibilis. Huic enim forma \((c,b,a)\) a parte prima erit contigua (art. 160) adeoque proprie aequivalens: sed \((c,b,a)\) per art. 159 formae \((a,b,c)\) improprie aequivalet: quare \((a,b,c)\) sibi ipsa improprie aequivalebit. Tales formas \((a,b,c)\) in quibus \(2b\) per \(a\) est divisibilis, \textit{ancipites} vocabimus. Habebimus itaque theorema hoc:

\textit{Forma \(F\), aliam formam \(F'\) tum proprie tum improprie implicabit, si forma anceps inveniri potest sub \(F\) contenta ipsam \(F'\) vero implicans.} Sed haec propositio etiam converti potest: scilicet

\Needspace{18\baselineskip}
\begin{center}
\textit{Theorema circa casum ubi forma sub alia simul proprie et improprie contenta est.}
\end{center}
\Article{164}
\textsc{Theorema.} \textit{Si forma}
\[
Axx+2Bxy+Cyy\quad\mathinner{.\,.\,.\,.\,.}\quad(F)
\]
\textit{formam}
\[
A'x'x'+2B'x'y'+C'y'y'\quad\mathinner{.\,.\,.\,.}\quad(F')
\]
\textit{tum proprie tum improprie implicat: forma anceps inveniri potest, sub \(F\) contenta formamque \(F'\) implicans.}

Ponamus, formam \(F\) transire in formam \(F'\) tum per substitutionem
\[
x=\alpha x'+\beta y',\qquad y=\gamma x'+\delta y'
\]
\SourceLeaf{141}
tum per hanc illi dissimilem
\[
x=\alpha'x'+\beta'x',\qquad y=\gamma'x'+\delta'y'
\]
Tum designatis numeris \(\alpha\delta-\beta\gamma,\ \alpha'\delta'-\beta'\gamma'\) per \(e,e'\), erit \(B'B'-A'C'=ee(BB-AC)=e'e'(BB-AC)\); hinc \(ee=e'e'\), et, quia per hyp. \(e,e'\) signa opposita habent, \(e=-e'\) sive \(e+e'=0\). Iam patet si in \(F'\) pro \(x'\) substituatur \(\delta'x''-\beta'y''\), et pro \(y',-\gamma'x''+\alpha'y''\), eandem formam esse prodituram ac si in \(F\) scribatur
\[
\begin{aligned}
\text{\textit{aut} 1) pro }x\quad&\alpha(\delta'x''-\beta'y'')+\beta(-\gamma'x''+\alpha'y'')\\
\text{\textit{i. e.}}\quad&(\alpha\delta'-\beta\gamma')x''+(\beta\alpha'-\alpha\beta')y''\\
\text{et pro }y\quad&\gamma(\delta'x''-\beta'y'')+\delta(-\gamma'x''+\alpha'y'')\\
\text{\textit{i. e.}}\quad&(\gamma\delta'-\delta\gamma')x''+(\delta\alpha'-\gamma\beta')y''\\
\text{\textit{aut} 2) pro }x\quad&\alpha'(\delta'x''-\beta'y'')+\beta'(-\gamma'x''+\alpha'y'')\quad\text{\textit{i. e.}}\quad e'x''\\
\text{et pro }y\quad&\gamma'(\delta'x''-\beta'y'')+\delta'(-\gamma'x''+\alpha'y'')\quad\text{\textit{i. e.}}\quad e'y''
\end{aligned}
\]
Designatis itaque numeris \(\alpha\delta'-\beta\gamma',\ \beta\alpha'-\alpha\beta',\ \gamma\delta'-\delta\gamma',\ \delta\alpha'-\gamma\beta'\) per \(a,b,c,d\): forma \(F\) per duas substitutiones
\[
x=ax''+by'',\quad y=cx''+dy'';\qquad x=e'x'',\quad y=e'y''
\]
in eandem formam transmutabitur, unde obtinemus tres aequationes sequentes:
\[
\begin{aligned}
Aaa+2Bac+Ccc&=Ae'e'&\quad\mathinner{.\,.\,.\,.\,.\,.\,.}\quad&[1]\\
Aab+B(ad+bc)+Ccd&=Be'e'&\quad\mathinner{.\,.\,.\,.\,.}\quad&[2]\\
Abb+2Bbd+Cdd&=Ce'e'&\quad\mathinner{.\,.\,.\,.\,.\,.\,.}\quad&[3]
\end{aligned}
\]
Ex valoribus ipsorum \(a,b,c,d\) autem invenitur
\[
ad-bc=ee'=-ee=-e'e'\quad\mathinner{.\,.\,.\,.\,.\,.}\quad[4]
\]
Hinc fit ex \(d[1]-c[2]\)
\[
(Aa+Bc)(ad-bc)=(Ad-Bc)e'e'
\]
adeoque
\[
A(a+d)=0
\]
Porro ex \((a+d)[2]-b[1]-c[3]\) fit
